\begin{table}[htbp]
\centering
\caption{Decomposition of the gender gap among high-achieving applicants and assignees, by field}
\label{tab:crossfield_decomposition}
\footnotesize
\setlength{\tabcolsep}{3pt}
\begin{tabular}{lccccccc}
\toprule
 & Pipeline & \multicolumn{3}{c}{Choice component $C$} & & Assignment & Balanced-pipeline \\
\cmidrule(lr){3-5}
Setting & $P$ & STEM & Health & Business/Law & & $V_{\text{STEM}}$ & male share $\Lambda(C_{\text{STEM}})$ \\
\midrule
Australia & $-$0.15 & +0.48 & $-$0.25 & +0.64 & & $-$0.05 & 0.62 \\
Brazil & +0.02 & +0.70 & $-$0.60 & +0.06 & & +0.05 & 0.67 \\
Chile & +0.24 & +0.77 & $-$0.67 & +0.12 & & $-$0.04 & 0.68 \\
Mainland China & $-$0.21 & +1.07 & $-$0.10 & $-$0.74 & & $-$0.17 & 0.75 \\
Finland & $-$0.07 & +0.78 & $-$0.52 & +0.40 & & $-$0.10 & 0.69 \\
Greece & $-$0.45 & +0.64 & +0.08 & -- & & +0.01 & 0.65 \\
Spain & $-$0.52 & +0.81 & $-$0.46 & $-$0.30 & & $-$0.08 & 0.69 \\
Sweden & $-$0.66 & +0.91 & $-$0.40 & $-$0.11 & & +0.06 & 0.71 \\
Taiwan & +0.43 & +0.37 & +0.04 & $-$0.83 & & +0.08 & 0.59 \\
Uganda & +0.46 & +0.29 & $-$0.06 & $-$0.88 & & +0.28 & 0.57 \\
\midrule
Positive in & 4/10 & 10/10 & 2/10 & 4/9 & & 5/10 & \\
\bottomrule
\end{tabular}
\begin{minipage}{\textwidth}\vspace{4pt}\scriptsize \textit{Notes:} High-achieving students are those in the top decile of each setting's academic performance distribution. $P$ is the pipeline component and $C$ the choice component of \eqref{eq:decomposition}; $V$ is the assignment component of \eqref{eq:decomposition_assign}, which extends the identity to realized assignments. Positive values favor men, and $P$ is common to all fields within a setting by construction. $\Lambda(C)$ is the male share among high-achieving STEM applicants that would obtain if the pipeline were balanced. Online \ref{sec:appendix_decomposition} extends the identity to assignments and reports $V$ for the other fields. Greece does not report business/law counts at the top decile. Rates here use each setting's full top-decile totals as denominators, so that the components satisfy the identity exactly; panel (b) of Figure \ref{fig:editor_fig2_pipeline} instead conditions Chile's rates on students who submit an application, so Chile's $C$ and its plotted gap are not directly comparable. \end{minipage}
\end{table}
